\documentclass[a4paper,12pt,twoside]{memoir}

% Castellano
\usepackage[spanish,es-tabla]{babel}
\selectlanguage{spanish}
\usepackage[utf8]{inputenc}
\usepackage[T1]{fontenc}
\usepackage{lmodern} % scalable font
\usepackage{microtype}
\usepackage{placeins}

\RequirePackage{booktabs}
\RequirePackage[table]{xcolor}
\RequirePackage{xtab}
\RequirePackage{multirow}

% Links
\usepackage[colorlinks]{hyperref}
\hypersetup{
	allcolors = {red}
}

% Ecuaciones
\usepackage{amsmath}

% Rutas de fichero / paquete
\newcommand{\ruta}[1]{{\sffamily #1}}

% Párrafos
\nonzeroparskip


% Imagenes
\usepackage{graphicx}
\newcommand{\imagen}[2]{
	\begin{figure}[!h]
		\centering
		\includegraphics[width=0.9\textwidth]{#1}
		\caption{#2}\label{fig:#1}
	\end{figure}
	\FloatBarrier
}

\newcommand{\imagenflotante}[2]{
	\begin{figure}%[!h]
		\centering
		\includegraphics[width=0.9\textwidth]{#1}
		\caption{#2}\label{fig:#1}
	\end{figure}
}



% El comando \figura nos permite insertar figuras comodamente, y utilizando
% siempre el mismo formato. Los parametros son:
% 1 -> Porcentaje del ancho de página que ocupará la figura (de 0 a 1)
% 2 --> Fichero de la imagen
% 3 --> Texto a pie de imagen
% 4 --> Etiqueta (label) para referencias
% 5 --> Opciones que queramos pasarle al \includegraphics
% 6 --> Opciones de posicionamiento a pasarle a \begin{figure}
\newcommand{\figuraConPosicion}[6]{%
  \setlength{\anchoFloat}{#1\textwidth}%
  \addtolength{\anchoFloat}{-4\fboxsep}%
  \setlength{\anchoFigura}{\anchoFloat}%
  \begin{figure}[#6]
    \begin{center}%
      \Ovalbox{%
        \begin{minipage}{\anchoFloat}%
          \begin{center}%
            \includegraphics[width=\anchoFigura,#5]{#2}%
            \caption{#3}%
            \label{#4}%
          \end{center}%
        \end{minipage}
      }%
    \end{center}%
  \end{figure}%
}

%
% Comando para incluir imágenes en formato apaisado (sin marco).
\newcommand{\figuraApaisadaSinMarco}[5]{%
  \begin{figure}%
    \begin{center}%
    \includegraphics[angle=90,height=#1\textheight,#5]{#2}%
    \caption{#3}%
    \label{#4}%
    \end{center}%
  \end{figure}%
}
% Para las tablas
\newcommand{\otoprule}{\midrule [\heavyrulewidth]}
%
% Nuevo comando para tablas pequeñas (menos de una página).
\newcommand{\tablaSmall}[5]{%
 \begin{table}
  \begin{center}
   \rowcolors {2}{gray!35}{}
   \begin{tabular}{#2}
    \toprule
    #4
    \otoprule
    #5
    \bottomrule
   \end{tabular}
   \caption{#1}
   \label{tabla:#3}
  \end{center}
 \end{table}
}

%
%Para el float H de tablaSmallSinColores
\usepackage{float}

%
% Nuevo comando para tablas pequeñas (menos de una página).
\newcommand{\tablaSmallSinColores}[5]{%
 \begin{table}[H]
  \begin{center}
   \begin{tabular}{#2}
    \toprule
    #4
    \otoprule
    #5
    \bottomrule
   \end{tabular}
   \caption{#1}
   \label{tabla:#3}
  \end{center}
 \end{table}
}

\newcommand{\tablaApaisadaSmall}[5]{%
\begin{landscape}
  \begin{table}
   \begin{center}
    \rowcolors {2}{gray!35}{}
    \begin{tabular}{#2}
     \toprule
     #4
     \otoprule
     #5
     \bottomrule
    \end{tabular}
    \caption{#1}
    \label{tabla:#3}
   \end{center}
  \end{table}
\end{landscape}
}

%
% Nuevo comando para tablas grandes con cabecera y filas alternas coloreadas en gris.
\newcommand{\tabla}[6]{%
  \begin{center}
    \tablefirsthead{
      \toprule
      #5
      \otoprule
    }
    \tablehead{
      \multicolumn{#3}{l}{\small\sl continúa desde la página anterior}\\
      \toprule
      #5
      \otoprule
    }
    \tabletail{
      \hline
      \multicolumn{#3}{r}{\small\sl continúa en la página siguiente}\\
    }
    \tablelasttail{
      \hline
    }
    \bottomcaption{#1}
    \rowcolors {2}{gray!35}{}
    \begin{xtabular}{#2}
      #6
      \bottomrule
    \end{xtabular}
    \label{tabla:#4}
  \end{center}
}

%
% Nuevo comando para tablas grandes con cabecera.
\newcommand{\tablaSinColores}[6]{%
  \begin{center}
    \tablefirsthead{
      \toprule
      #5
      \otoprule
    }
    \tablehead{
      \multicolumn{#3}{l}{\small\sl continúa desde la página anterior}\\
      \toprule
      #5
      \otoprule
    }
    \tabletail{
      \hline
      \multicolumn{#3}{r}{\small\sl continúa en la página siguiente}\\
    }
    \tablelasttail{
      \hline
    }
    \bottomcaption{#1}
    \begin{xtabular}{#2}
      #6
      \bottomrule
    \end{xtabular}
    \label{tabla:#4}
  \end{center}
}

%
% Nuevo comando para tablas grandes sin cabecera.
\newcommand{\tablaSinCabecera}[5]{%
  \begin{center}
    \tablefirsthead{
      \toprule
    }
    \tablehead{
      \multicolumn{#3}{l}{\small\sl continúa desde la página anterior}\\
      \hline
    }
    \tabletail{
      \hline
      \multicolumn{#3}{r}{\small\sl continúa en la página siguiente}\\
    }
    \tablelasttail{
      \hline
    }
    \bottomcaption{#1}
  \begin{xtabular}{#2}
    #5
   \bottomrule
  \end{xtabular}
  \label{tabla:#4}
  \end{center}
}



\definecolor{cgoLight}{HTML}{EEEEEE}
\definecolor{cgoExtralight}{HTML}{FFFFFF}

%
% Nuevo comando para tablas grandes sin cabecera.
\newcommand{\tablaSinCabeceraConBandas}[5]{%
  \begin{center}
    \tablefirsthead{
      \toprule
    }
    \tablehead{
      \multicolumn{#3}{l}{\small\sl continúa desde la página anterior}\\
      \hline
    }
    \tabletail{
      \hline
      \multicolumn{#3}{r}{\small\sl continúa en la página siguiente}\\
    }
    \tablelasttail{
      \hline
    }
    \bottomcaption{#1}
    \rowcolors[]{1}{cgoExtralight}{cgoLight}

  \begin{xtabular}{#2}
    #5
   \bottomrule
  \end{xtabular}
  \label{tabla:#4}
  \end{center}
}




\graphicspath{ {./img/} }

% Capítulos
\chapterstyle{bianchi}
\newcommand{\capitulo}[2]{
	\setcounter{chapter}{#1}
	\setcounter{section}{0}
	\chapter*{#2}
	\addcontentsline{toc}{chapter}{#2}
	\markboth{#2}{#2}
}

% Apéndices
\renewcommand{\appendixname}{Apéndice}
\renewcommand*\cftappendixname{\appendixname}

\newcommand{\apendice}[1]{
	%\renewcommand{\thechapter}{A}
	\chapter{#1}
}

\renewcommand*\cftappendixname{\appendixname\ }

% Formato de portada
\makeatletter
\usepackage{xcolor}
\newcommand{\tutor}[1]{\def\@tutor{#1}}
\newcommand{\course}[1]{\def\@course{#1}}
\definecolor{cpardoBox}{HTML}{E6E6FF}
\def\maketitle{
  \null
  \thispagestyle{empty}
  % Cabecera ----------------
\noindent\includegraphics[width=\textwidth]{cabecera}\vspace{1cm}%
  \vfill
  % Título proyecto y escudo informática ----------------
  \colorbox{cpardoBox}{%
    \begin{minipage}{.8\textwidth}
      \vspace{.5cm}\Large
      \begin{center}
      \textbf{TFG del Grado en Ingeniería Informática}\vspace{.6cm}\\
      \textbf{\LARGE\@title{}}
      \end{center}
      \vspace{.2cm}
    \end{minipage}

  }%
  \hfill\begin{minipage}{.20\textwidth}
    \includegraphics[width=\textwidth]{escudoInfor}
  \end{minipage}
  \vfill
  % Datos de alumno, curso y tutores ------------------
  \begin{center}%
  {%
    \noindent\LARGE
    Presentado por \@author{}\\ 
    en Universidad de Burgos --- \@date{}\\
    Tutor: \@tutor{}\\
  }%
  \end{center}%
  \null
  \cleardoublepage
  }
\makeatother


% Datos de portada
\title{título del TFG \\Documentación Técnica}
\author{nombre alumno}
\tutor{nombre tutor}
\date{\today}

\begin{document}

\maketitle



\cleardoublepage



%%%%%%%%%%%%%%%%%%%%%%%%%%%%%%%%%%%%%%%%%%%%%%%%%%%%%%%%%%%%%%%%%%%%%%%%%%%%%%%%%%%%%%%%



\frontmatter


\clearpage

% Indices
\tableofcontents

\clearpage

\listoffigures

\clearpage

\listoftables

\clearpage

\mainmatter

\appendix

\apendice{Plan de Proyecto Software}

\section{Introducción}

En este documento hablaremos de la planificación del proyecto, que es una parte muy importante. Aquí analizaremos el tiempo de desarrollo y el presupuesto que se necesitara para llevarlo a cabo.

lo dividiremos en dos partes:

\begin{itemize}
	\tightlist
	\item
	Planificación temporal del proyecto.
	\item
	Estudio de viabilidad del proyecto.
\end{itemize}

En la planificación temporal hablaremos de cada uno de los \emph{sprints} necesarios para su realización, el tiempo estimado y el tiempo real de realización.

En el estudio de viabilidad hablaremos de los costes y beneficios que nos aporta esta aplicación y su viabilidad legal, que se refiere a las licencias de uso del código ajeno que hemos utilizado en el desarrollo para poder comercializar nuestra aplicación.

\section{Planificación temporal}

Al inicio del proyecto se propuso utilizar una metodología ágil, en concreto Scrum, ya que se decidió tener reuniones semanales para hablar de los cambios realizados, problemas ocasionados y planificación del siguiente Sprint que se realizara. No se ha conseguido desarrollar la metodología al 100\% ya que el el equipo de desarrollo constaba de 1 persona (Oscar Fernández Armengol), pero desechando este punto, se ha conseguido un desarrollo ágil en sus demás puntos. 

\begin{itemize}
	\tightlist
	\item
	Se aplicó una estrategia de desarrollo incremental a través de
	iteraciones (\emph{sprints}) y revisiones.
	\item
	La duración media de los \emph{sprints} fue de una semana.
	\item
	Al finalizar cada \emph{sprint} se entregaba un producto funcional con la nueva especificación en el caso de que estuviera terminada.
	\item
	Se realizaban reuniones de revisión al finalizar cada \emph{sprint}, de resolución de dudas y
	al mismo tiempo de planificación del nuevo \emph{sprint}.
	\item
	En la planificación del \emph{sprint} se generaba una lista de tareas a realizar (nuevas funcionalidades o bugs a solucionar).
	\item
	Se estimaba el tiempo de realización de las tareas a realizar en el \emph{canvas}.
	\item
	Para monitorizar el progreso del proyecto se utilizan los gráficos generados en github.
	
\end{itemize}

\subsection{Sprint 1 (03/10/17 -
	10/10/17)}\label{sprint-1-0031017---101017}
Este \emph{sprint} fue el comienzo del proyecto, aunque en reuniones previas se hablo con el tutor de las propuestas que tenia para la elección del proyecto, una vez el tutor (Raúl Marticorena Sanchez) acepto tutorizar al alumno (Oscar Fernández Armengol) se puedo empezar el desarrollo.

Los objetivos fueron: preparación del entorno de desarrollo
\begin{itemize}
	\tightlist
	\item
	Preparación del entorno de desarrollo.
	\item
	Familiarización con la aplicación heredada \href{https://github.com/claumartinezh/TFG_UBUGrades}{UBUGrades}.
	\item
	Investigación del web service de moodle
	\item
	Creación de un esqueleto del proyecto para poder empezar a trabajar.
	
\end{itemize}
El \href{https://github.com/trona85/GII-17.1B-UBULog-1.0/milestone/1?closed=1}{Sprint 1} se estimo en 6 días de trabajo y se realizo en esos 6 días.

\imagen{sprint1}{Sprint 1.}

\subsection{Sprint 2 (10/10/17 -
	17/10/17)}\label{sprint-2-101017---171017}

El \href{https://github.com/trona85/GII-17.1B-UBULog-1.0/milestone/2?closed=1}{Sprint 2}

\imagen{sprint2}{Sprint 2.}

\subsection{Sprint 3 (17/10/17 -
	24/10/17)}\label{sprint-3-171017---241017}

El \href{https://github.com/trona85/GII-17.1B-UBULog-1.0/milestone/3?closed=1}{Sprint 3} 

\imagen{sprint3}{Sprint 3.}

\subsection{Sprint 4 (24/10/17 -
	7/11/17)}\label{sprint-4-241017---071117}

El \href{https://github.com/trona85/GII-17.1B-UBULog-1.0/milestone/4?closed=1}{Sprint 4} 

\imagen{sprint4}{Sprint 4.}

\subsection{Sprint 5 (7/11/17 -
	14/11/17)}\label{sprint-5-071117---141117}

El \href{https://github.com/trona85/GII-17.1B-UBULog-1.0/milestone/5?closed=1}{Sprint 5} 

\imagen{sprint5}{Sprint 5.}

\subsection{Sprint 6 (14/11/17 -
	22/11/17)}\label{sprint-6-141117---221117}

El \href{https://github.com/trona85/GII-17.1B-UBULog-1.0/milestone/6?closed=1}{Sprint 6} 

\imagen{sprint6}{Sprint 6.}

\subsection{Sprint 7 (22/11/17 -
	29/11/17)}\label{sprint-7-221117---291117}

El \href{https://github.com/trona85/GII-17.1B-UBULog-1.0/milestone/7?closed=1}{Sprint 7} 

\imagen{sprint7}{Sprint 7.}

\subsection{Sprint 8 (28/11/17 -
	06/12/17)}\label{sprint-8-281117---061217}

El \href{https://github.com/trona85/GII-17.1B-UBULog-1.0/milestone/8?closed=1}{Sprint 8} 

\imagen{sprint8}{Sprint 8.}

\subsection{Sprint 9 (06/12/17 -
	13/12/17)}\label{sprint-8-061217---131217}

El \href{https://github.com/trona85/GII-17.1B-UBULog-1.0/milestone/9?closed=1}{Sprint 9}




\section{Estudio de viabilidad}

\subsection{Viabilidad económica}

\subsection{Viabilidad legal}



\apendice{Especificación de Requisitos}

\section{Introducción}

Este proyecto está orientado a profesores de la universidad de Burgos y en este apartado vamos a hablar de los objetivos generales, de la especificación de requisitos (funcionales y no funcionales) y especificaremos cada uno de ellos.

\section{Objetivos generales}

En el desarrollo del proyecto se quieren conseguir los siguientes objetivos:

\begin{itemize}
	\tightlist
	\item
	Desarrollar una aplicación de escritorio para poder loguearse con UBUVirtual.
	\item
	Obtención de datos por el Web Service de Moodle.
	\item
	Visualizar los datos correspondientes a la asignatura seleccionada.
	\item
	Facilitar la interpretación de los datos representándolo por diferentes tipos de gráficas.
	\item
	Desgranar los datos aportados para un análisis más profundo de ellos.
\end{itemize}

\section{Catálogo de requisitos}

Ahora hablaremos de los requisitos funcionales y no funcionales, derivados de los objetivos generales de este proyecto. Se notará que varios puntos del apartado son iguales al anexo del TFG de Claudia Martínez Herrero \cite{claudia:anexo}.

\subsection{Requisitos funcionales}\label{requisitos-funcionales}

\begin{itemize}
	\tightlist
	\item
	\textbf{RF-1 Autenticación de usuario:} Inicialmente el usuario debe identificarse con su correo de profesor y su contraseña para poder hacer uso de las funciones de la aplicación.
	\item
	\textbf{RF-2 Extracción de datos:} A partir de los datos de usuario, se obtendrán todos los datos referidos a los cursos en los que esté matriculado.
	\item
	\textbf{RF-3 Carga de datos:} El usuario debe poder cargar los datos correspondientes al log de la asignatura.
	\item
	\textbf{RF-4 Selección de datos:} El usuario debe poder seleccionar qué participantes de esa asignatura y en qué eventos quiere que se muestre el gráfico.
	
	\begin{itemize}
		\tightlist
		\item
		\textbf{RF-4.1 Filtrado de participantes y eventos:} Se podrá filtrar a los participantes según los grupos, el rol y por un campo de texto y los eventos por un campo de texto.
	\end{itemize}
	\item
	\textbf{RF-5 Visualización de gráficas de logs:} El usuario podrá visualizar los registros en forma de gráfica.
	\item
	\textbf{RF-6 Visualización de Logs:} El usuario podrá visualizar las registros de manera gráfica en formato de tabla.
		\begin{itemize}
			\tightlist
			\item
			\textbf{RF-6.1 Filtrado de logs:} Se podrá filtrar logs por varios campos de texto cada uno asociado a una columna de la tabla.
		\end{itemize}
	\item
	\textbf{RF-7 Cierre de sesión:} La aplicación cierra sesión.
\end{itemize}

\subsection{Requisitos no funcionales}\label{requisitos-no-funcionales}

\begin{itemize}
	\tightlist
	\item
	\textbf{RNF-1} La aplicación debe ser de escritorio.
	\item
	\textbf{RNF-2} La aplicación debe funcionar en un ordenador con Java (versión 8)
	\item
	\textbf{RNF-3} La aplicación debe funcionar contra plataformas de Moodle 3.3 como es UBUVirtual.
	\item
	\textbf{RNF-4} La interfaz de usuario debe ser clara e intuitiva para el usuario.
	\item
	\textbf{RNF-5} La aplicación debe ser sencilla de utilizar, ahorrando pasos innecesarios que puedan resultar confusos para el usuario.
	\item
	\textbf{RNF-6} La aplicación será extensible.
	\item
	\textbf{RNF-7} La sección de gráficos debe permitir exportar los gráficos en formato de imagen: .png, .jpg....

\end{itemize}

\section{Especificación de requisitos}

En esta sección se mostrará el diagrama de casos de uso resultante y se
desarrollará cada uno de ellos.	
\subsection{Diagrama de casos de uso}\label{diagrama-de-casos-de-uso}
	
\imagen{usecasesdiagram}{Diagrama de casos de uso.}{1.35}

\subsection{Actores}\label{actores}

Solo interactuará con el sistema un actor, que se corresponderá con la
figura del usuario.

\subsection{Casos de uso}\label{casos-de-uso}

\tablaSmallSinColores{RF1 - Autenticación de usuario}{|l | l |}{autenticaciondeusuario}
{ \multicolumn{1}{|l }{RF1} & Autenticación de usuario \\}{ 
	Descripción & Se introduce e-mail, contraseña y host (UBUVirtual).\\ \hline
	Requisitos asociados & RF2. \\
	\hline
	Precondiciones & El usuario debe tener utilizar el e-mail de UBUVirtual. \\ \hline
	Secuencia normal & 	 1- Introducir los 3 campos del login. \\
	
	 & 2- Entrar en la aplicación. \\ \hline
	Postcondiciones & Datos de usuario cargados. Se accederá a la pantalla de \\
	 & bienvenida.\\ \hline
	Frecuencia & Alta.\\ \hline
	Importancia & Alta. \\ 
} 

\tablaSmallSinColores{RF2 - Extracción de datos}{|l | l |}{extraccion-de-datos}
{ \multicolumn{1}{|l }{RF2} & Extracción de datos \\}{ 
	Descripción & El sistema carga los cursos en los que está matriculado \\  & el usuario.\\ \hline
	Requisitos asociados & RF1. \\
	\hline
	Precondiciones &El usuario se ha logueado correctamente. \\ \hline
	Secuencia normal & 	 1- Obtención de los cursos del usuario. \\
	& 2- Carga de la lista de cursos. \\
		& 3- Selección de curso por parte del usuario. \\ \hline
	Postcondiciones & Se visualiza una nueva pantalla con el curso elegido.\\ \hline
	Frecuencia & Alta.\\ \hline
	Importancia & Alta. \\ 
}

\tablaSmallSinColores{RF3 - Carga de datos}{|l | l |}{carga-de-datos}
{ \multicolumn{1}{|l }{RF3} & Extracción de datos \\}{ 
	Descripción & El usuario cargara el fichero csv o ejecutara la descarga del log.\\ \hline
	Requisitos asociados & RF1. \\
	\hline
	Precondiciones & 1- El usuario se ha logueado correctamente. \\ 
	& 2- El usuario selecciona una asignatura. \\ \hline
	Secuencia normal & 	 SI tiene csv: Carga documento csv. \\
	& SI NO tiene csv: puede descargarlo manualmente o clicando en\\ & descarga automatizada. \\ \hline
	Postcondiciones & Se cargan los datos de log en la aplicación.\\ \hline
	Frecuencia & Media.\\ \hline
	Importancia & Alta. \\ 
}

\tablaSmallSinColores{RF4 - Selección de datos}{|l | l |}{selección-de-datos}
{ \multicolumn{1}{|l }{RF4} & Selección de datos \\}{ 
	Descripción & Se visualiza una pantalla con la lista de participantes del curso\\
	& y la lista de eventos. El usuario debe seleccionar qué participantes \\ &que eventos o ambos desea visualizar.\\ \hline
	Requisitos asociados & RF3, RF4.1. \\
	\hline
	Precondiciones & 1- El log debe de estar cargado. \\ 
	& 2- Se debe mostrar participantes y eventos disponibles. \\
	& 3- Se debe clicar en participantes o eventos o ambos. \\ \hline
	Secuencia normal & 1- Se carga el log. \\ 
	& 2- Se muestran los participantes y eventos disponibles. \\
	& 3- Se hace clic en participantes o eventos o ambos. \\ \hline
	Postcondiciones & Se cargan los datos de log en la aplicación.\\ \hline
	Frecuencia & Alta.\\ \hline
	Importancia & Alta. \\ 
}

\tablaSmallSinColores{RF4.1 - Filtrado de participantes y eventos}{|l | l |}{filtrado-de-participantes-y-eventos}
{ \multicolumn{1}{|l }{RF4} & Filtrado de participantes y eventos \\}{ 
	Descripción & El usuario podrá filtrar los participantes y eventos del\\ & curso.\\ \hline
	Requisitos asociados & RF4. \\
	\hline
	Precondiciones & 1- El log debe de estar cargado. \\ 
	& 2- Se debe mostrar participantes y eventos disponibles. \\ \hline
	Secuencia normal & 1- Se carga el log. \\ 
	& 2- Se muestran los participantes y eventos disponibles. \\
	& 3- Se puede interaccionar con los filtros. \\ \hline
	Postcondiciones & Se muestran los datos que coinciden con los filtros.\\ \hline
	Frecuencia & Alta.\\ \hline
	Importancia & Alta. \\ 
}

\tablaSmallSinColores{RF5 - Visualización de gráficas de logs}{|l | l |}{visualización-de-gráficas-de-logs}
{ \multicolumn{1}{|l }{RF5} & Visualización de gráficas de logs \\}{ 
	Descripción & La aplicación mostrara los datos seleccionados de manera\\ & gráfica.\\ \hline
	Requisitos asociados & RF3, RF4. \\
	\hline
	Precondiciones & 1- El log debe de estar cargado. \\ 
	& 2- Se debe mostrar participantes y eventos disponibles. \\
	& 3- Se hace clic en participantes o eventos o ambos. \\ \hline
	Secuencia normal & 1- Se carga el log. \\ 
	& 2- Se muestran los participantes y eventos disponibles. \\
	& 3- Se hace clic en participantes o eventos o ambos. \\ \hline
	Postcondiciones & Se muestran los datos en la gráfica.\\ \hline
	Frecuencia & Alta.\\ \hline
	Importancia & Alta. \\ 
}

\tablaSmallSinColores{RF6 - Visualización de logs}{|l | l |}{visualización-de-logs}
{ \multicolumn{1}{|l }{RF6} & Visualización de logs \\}{ 
	Descripción & La aplicación mostrara los logs cargados en forma de tabla.\\ \hline
	Requisitos asociados & RF3, RF4. \\
	\hline
	Precondiciones & 1- El log debe de estar cargado. \\ 
	& 2- Se debe mostrar participantes y eventos disponibles. \\
	& 3- Se hace clic en participantes o eventos o ambos. \\ \hline
	Secuencia normal & 1- Se carga el log. \\ 
	& 2- Se muestran los participantes y eventos disponibles. \\
	& 3- Se hace clic en participantes o eventos o ambos. \\ \hline
	Postcondiciones & Se muestran los datos en la tabla logs.\\ \hline
	Frecuencia & Alta.\\ \hline
	Importancia & Alta. \\ 
}

\tablaSmallSinColores{RF6.1 - Filtrado de logs}{|l | l |}{filtrado-de-logs}
{ \multicolumn{1}{|l }{RF6.1} & Filtrado de logs \\}{ 
	Descripción & El usuario podrá filtrar logs del curso.\\ \hline
	Requisitos asociados & RF6. \\
	\hline
	Precondiciones & 1- El log debe de estar cargado. \\ 
	& 2- Se debe mostrar participantes y eventos disponibles. \\ \hline
	Secuencia normal & 1- Se carga el log. \\ 
	& 2- Se muestran los participantes y eventos disponibles. \\
	& 3- Se puede interaccionar con los filtros. \\ \hline
	Postcondiciones & Se muestran los datos que coinciden con los filtros.\\ \hline
	Frecuencia & Alta.\\ \hline
	Importancia & Alta. \\ 
}

\tablaSmallSinColores{RF7 - Cierre de sesión de usuario}{|l | l |}{cierre-de-sesión-de-usuario}
{ \multicolumn{1}{|l }{RF7} & Cierre de sesión de usuario \\}{ 
	Descripción & Mediante un botón de salir el usuario podrá cerrar la\\ & aplicación.\\ \hline
	Requisitos asociados & Ninguno. \\
	\hline
	Precondiciones & Sesión de usuario activa. \\ \hline
	Secuencia normal & 1- El usuario pulsa el botón de salir. \\ 
	& 2- El sistema cierra la aplicación. \\ \hline
	Postcondiciones & Aplicación cerrada.\\ \hline
	Frecuencia & Media.\\ \hline
	Importancia & Media. \\ 
}





\apendice{Especificación de diseño}

\section{Introducción}

\section{Diseño de datos}

\section{Diseño procedimental}

\section{Diseño arquitectónico}



\apendice{Documentación técnica de programación}

\section{Introducción}

En este anexo describe la documentación técnica de programación,
incluyendo la instalación del entorno de desarrollo, la estructura de la
aplicación, su compilación, ejecución, la configuración de los diferentes servicios y las pruebas del sistema.

\section{Estructura de directorios}

El repositorio del proyecto podemos observar la siguiente estructura:

\begin{itemize}
	\tightlist
	\item
	\texttt{/}: contiene los ficheros de configuración de Gradle \cite{gradle}, de los servicios de integración continua, en nuestro caso Travis-CI \cite{web:travis}, el fichero README , la copia de la
	licencia y un fichero .xml que es el ANT para la creación del .jar.
	\item
	\texttt{/src/main/java/}: Encontraremos toda la lógica de la aplicación.
	\item
	\texttt{/src/test/java/}: El código para testear la aplicación.
	\item
	\texttt{/resources/}: Recursos de nuestra aplicación.
	\item
	\texttt{/resources/css/}: Estilos de la aplicación.
	\item
	\texttt{/resources/img/}: Imágenes de la aplicación.
	\item
	\texttt{/lib/}: Librerías de la aplicación.
	\item
	\texttt{/doc/}: documentación del proyecto.
	\item
	\texttt{/doc/javadoc/}: documentación \emph{javadoc}.
	\item
	\texttt{/docs/proyectdoc/}: documentación en formato \LaTeX.
\end{itemize}

\section{Manual del programador}

En esta sección explicaremos como montar el entorno de desarrollo, descargar el código fuente del proyecto, compilarlo, ejecutarlo y exportarlo. hay que tener en cuenta que este manual esta orientado a sistemas Linux como lo es Ubuntu 16.04 que es donde se ha desarrollado el proyecto, aunque también se podría desarrollar en S.O. Windows.

Para la instalación en sistemas Windows hay que leer la documentación de Claudia Martínez Herrero \cite{claudia:anexo} referente al manual del programador.

\subsection{Entorno de desarrollo}\label{entorno-de-desarrollo}

Para trabajar con el proyecto se necesita tener instalados los
siguientes programas y dependencias:

\begin{itemize}
	\tightlist
	\item
	Java JDK 8.
	\item
	Git.
	\item
	Moodle 3.3.
	\item
	Eclipse.
	\item
	Plugin Gradle para Eclipse.
\end{itemize}

A continuación, se indica como instalar y configurar correctamente cada
uno de ellos.

\subsubsection{Java JDK 8}\label{java-jdk-8}

En esta instalación es importante descargar el JDK 8 de oracle ya que, otro jdk como puede ser el que viene por defecto en las librerías de Ubuntu no te proporcionan la Librería Java FX.

Deberemos abrir la consola y poner los siguientes comandos:

\imagen{jdk}{Instalacion JDK 8 en Linux.}

\subsubsection{Git}\label{git}

Necesitamos git para la gestión del repositorio. Git nos permitirá clonar el proyecto en nuestro equipo, subir cambios, comprobar diferencias entre commit, etc.

\imagen{git}{Instalacion Git.}

\subsubsection{Moodle 3.3}\label{moodle}

Para la instalación de Moodle hacemos referencia a los anexos de Claudia Martínez Herrero \cite{claudia:anexo} referente a la instalación de moodle en Windows ya que es idéntica.

\subsubsection{Eclipse}\label{git}

Eclipse es el entorno de desarrollo seleccionado para este proyecto. Si el programador desea utilizar otro, no tendría ningún problema.

Para la instalación deberemos descargar \href{http://www.eclipse.org/downloads/}{Eclipse} en nuestro equipo. Lo descomprimimos y ejecutamos el archivo eclipse-inst, los pasos siguientes es igual que en Windows, se abrirá el gestor de instalación y solo hay que seguir las instrucciones.

\subsubsection{Gadle}\label{git}

La mayoría de los IDE de desarrollo tienen integrado Gradle, como es el caso de Eclipse y no se necesita ningún tipo de instalación.

En el caso de que no tenga, podemos instalarlo en nuestra maquina personal. Lo que deberemos hacer es:

\imagen{gradle}{Instalación gradle.}

\section{Compilación, instalación y ejecución del proyecto}

\subsubsection{Obtención del código fuente}\label{obtención-del-código-fuente}

Una vez hemos instalado todo lo necesario, deberemos clonar nuestro proyecto.

Deberemos entrar al repositorio \href{https://github.com/trona85/GII-17.1B-UBULog-1.0}{GitHub} y en la parte que pone Clone whith HTTPS, copiarnos la URL.

\imagen{githubclone}{Clonar proyecto con Git 1.}

Abrimos la consola y nos vamos a nuestro workspace. Una vez ahí, ponemos el siguiente comando de git con la URL copiada anteriormente.

\imagen{gitclone}{Clonar proyecto con Git 2.}

Nos descargara el código de la aplicación y ya lo tendremos en nuestro equipo para poder trabajar con el.

\subsubsection{Importar proyecto en Eclipse}\label{importar-proyecto-en-eclipse}

Para importar nuestro proyecto deberemos abrir Eclipse. Una vez abierto deberemos:

\begin{enumerate}
	\tightlist
	\item
	Vamos a File-->import...
	\item
	Nos saldrá una ventana en la que deberemos elegir la opción de gradle.
	
	\imagen{importgradle}{Importar proyecto Gradle.}
	
	\item
	Nos saldra una ventana de información, después de su lectura se da a next.
	
		\imagen{welcomegradle}{Información Gradle.}
	\item
	En esta ventana nos ubicamos en la raíz del proyecto que queremos importar, igual que en la siguiente imagen, después pinchamos en finalizar.
	
	\imagen{raizgradle}{Cargar proyecto Gradle.}
	
	\item
	Gradle empezara a descargar todas las dependencias del proyecto, esto puede tardar unos minutos. Cuando termine ya tendremos nuestro proyecto preparado para el desarrollo.
		
	\imagen{dependeciesgradle}{Cargar proyecto Gradle.}
	
	
\end{enumerate}

\subsubsection{Build.gradle}\label{buildgradle}

Build.gradle es el archivo donde debemos poner el tipo de aplicación que es, la versión y las dependencias necesarias de nuestro proyecto. \cite{gradle}

\subsubsection{Travis.yml}\label{buildgradle}

Travis.yml es el archivo donde tenemos la configuración para la integración continua y el análisis de la calidad del código con sonarqube. \cite{web:travis} \cite{Web:sonarqube}

\subsubsection{Generar .jar}\label{generar-jar}

Para generar el .jar se ha creado un .xml en el proyecto que es un ANT, en el se tienen todas las dependencias, .class y recursos necesarios para la ejecución. Para generar el jar deberemos:

\begin{enumerate}
	\tightlist
	\item
	Abrimos UBULog.xml en eclipse.
	\item
	La ventana de Eclipse que esta en la parte derecha, que pone OutLine tendremos una opción que es create runnable.
	
	\imagen{antecl}{Crear .jar.}
	
	Con el botón secundario pinchamos y ejecutamos el ANT Builder para generar el jar, en la consola del Eclipse nos saldrá esto.
	
	\imagen{antcreate}{Jar creado.}
	
	Una vez visto el resultado de la captura anterior, querrá decir que ya hemos generado el .jar.
	
\end{enumerate}

%\section{Pruebas del sistema}

\apendice{Documentación de usuario}

\section{Introducción}

En este manual detallamos los requisitos necesarios y el manejo de la aplicación, para ayudar al usuario a su comprensión. No es necesaria la instalación, ya que es un ejecutable.

\section{Requisitos de usuarios}

Los requisitos mínimos para poder hacer uso de la aplicación son:

\begin{itemize}
	\tightlist
	\item
	Disponer del ejecutable de la aplicación, UBULog.jar.
	\item
	Disponer de un ordenador con conexión a Internet.
	\item
	Tener instalado Java 8 en el equipo que se vaya a ejecutar la aplicación.
	\item
	El usuario debe tener una cuenta de UBUVirtual activa con el rol de profesor.
\end{itemize}

%\section{Instalación}

\section{Manual del usuario}

\subsection{Ejecución}

Deberemos hacer doble clic en el ejecutable. En el caso de que no se nos habrá, será por qué no tenemos por defecto abrirlo con java 8, pinchando encima con el botón secundario podemos darle a ``abrir con'' y seleccionar java 8.

En el caso que queramos ejecutarlo desde la consola deberemos abrirla, movernos a la ubicación del ejecutable y poner ``java -jar nombreaplicación.jar''. En la siguiente captura mostramos el proceso.

\imagen{execut}{Ejecución UBULog.}{0.9}

IMPORTANTE: Se ha detectado que la librería Java FX da errores con algunas versiones de Java 8 cuando intenta mostrar las tablas de log y el gráfico, hay que tener instalado las ultimas versiones.

\subsection{Descargar el log manualmente}

Para descargar el log manualmente deberemos ir a UBUVirtual, iniciar sesión y irnos a la asignatura deseada. Una vez allí en la parte de administración de la asignatura deberemos pinchar en ``Registros''.

\imagen{administracion}{Administración de la asignatura.}{0.6}

Una vez pinchemos en registro, nos llevará a la siguiente ventana:

\imagen{filtross}{Filtros log.}{0.9}

Ponemos las opciones que deseemos y pinchamos en ``Conseguir estos registros''.

\imagen{filtrosscons}{Registros.}{0.9}

Despues en la parte inferior tendremos un botón que nos pondrá ``Descargar''.

\imagen{logsdescc}{Descarga de logs.}{0.9}

Una vez descargado ya lo podremos utilizar en nuestra aplicación.


\subsection{Login}


Una vez ejecutada la aplicación nos saldrá la ventana para iniciar sesión en nuestra aplicación.

\imagen{login}{Login UBULog.}{0.6}

En esta ventana observamos tres campos que deberemos rellenar:

\begin{itemize}
	\tightlist
	\item
	\textbf{Usuario:} Correo de la personal UBU.
	\item
	\textbf{Contraseña:} Contraseña utilizada para iniciar sesión en UBUVirtual.
	\item
	\textbf{Host:} Dirección web de la plataforma a donde nos queremos conectar. Por defecto está https://ubuvirtual.ubu.es/.
\end{itemize}

Una vez introduzcamos los campos tenemos dos opciones:

\begin{itemize}
	\tightlist
	\item
	\textbf{Clicar borrar:} Se borrarán los campos.
	\item
	\textbf{Clicar entrar:} Pueden suceder 2 cosas:
	\begin{itemize}
		\tightlist
		\item
		\textbf{Datos incorrectos:} Se borrarán los datos e indicará al usuario que el usuario o la contraseña es errónea.
		\item
		\textbf{Datos correctos:} Se nos abrirá la pantalla de bienvenida.
	\end{itemize}
\end{itemize}

\subsection{Pantalla de bienvenida}


Cuando iniciemos sesión en nuestra aplicación llegaremos a la pantalla de bienvenida, en ella, nos aparecerá la lista de asignaturas en las que estamos matriculado. Deberemos seleccionar una en la que estemos matriculados con el rol de profesor.

\imagen{bienvenida}{Pantalla de bienvenida.}{0.9}

Cuando hagamos clic en entrar pueden pasar dos cosas:

\begin{itemize}
	\tightlist
	\item
	\textbf{Si NO hemos seleccionado asignatura:} La aplicación nos informa que debemos seleccionar una asignatura.
	\item
	\textbf{Si hemos seleccionado asignatura:} Nos llevará a la ventana principal de la aplicación.
	
\end{itemize}

\subsection{Pantalla principal}

\imagen{ubulog}{Pantalla principal.}{0.9}

\subsubsection{Visualización de participantes}

En esta parte de la aplicación veremos los participantes de la asignatura (matriculados, desconocidos, sistema y administrador), nos dirá el número de participantes y tendremos tres filtros:


\begin{itemize}
	\tightlist
	\item
	\textbf{Filtro participantes:} Es un filtro de tipo texto que filtrará por el nombre de los participantes.
	\item
	\textbf{Grupos:} Filtrará por los grupos a los que pertenece los participantes.
	\item
	\textbf{Rol:} Filtrará por los roles a los que pertenece los participantes.
	
\end{itemize}


\imagen{participantes}{Visualización de participantes.}{0.6}

Esta sección de la aplicación estará deshabilitada hasta que se cargue el log. Una vez cargado podremos interactuar con él para mostrar los datos en la gráfica y en la tabla de log.

\subsubsection{Visualización de eventos}

En esta parte de la aplicación veremos los eventos que se han realizado a lo largo de la asignatura y un filtro de texto para filtrar por el nombre del evento.


\imagen{evento}{Visualización de eventos.}{0.6}

Esta sección de la aplicación estará deshabilitada hasta que se cargue el log. Una vez cargado podremos interactuar con él para mostrar los datos en la gráfica y en la tabla de log.

\subsubsection{Gráfica}

En esta parte de la aplicación podremos ver los datos seleccionados de forma gráfica, en esta ventana, también tenemos la posibilidad de cambiar el tipo de gráfico.

\imagen{grafica}{Gráfica.}{0.9}

\subsubsection{Tabla logs}

En esta parte de la aplicación podremos ver los log correspondientes a los datos seleccionados en forma de tabla, en esta ventana, también tenemos la posibilidad de filtrar los datos con los filtros que se encuentra en la parte superior, cada uno corresponde a una columna de la tabla y no diferencian entre minúsculas y mayúsculas.

\imagen{tablelog}{Tabla logs.}{0.9}

\subsubsection{Botones inferiores}

En esta parte de la aplicación podremos clicar los siguientes botones y su comportamiento es el siguiente:

\begin{itemize}
	\tightlist
	\item
	\textbf{Generar gráfica a partir de filtros de tabla:} Este botón permanecerá deshabilitado hasta que se cargue el log. De lo que se encarga es de coger los log de la tabla resultante sin filtrar o si se han aplicado filtros y generar un nuevo gráfico con los nuevos datos.
	\item
	\textbf{Exportar gráfico:} Guarda el gráfico en una imagen.
	\item
	\textbf{Cargar documento online:} Carga los datos del log de forma automatizada. Este proceso puede tardar varios minutos.
	\item
	\textbf{Cargar documento local:} Habiendo descargado el log de la asignatura en CSV, podemos pasárselo a la aplicación directamente y cargar los datos. Este proceso es más eficiente que el punto anterior.
	\item
	\textbf{Salir:} Cierra la aplicación.
	
	
\end{itemize}

\imagen{botonesinferiores}{Botones inferiores.}{0.9}

Esta sección de la aplicación estará deshabilitada hasta que se cargue el log. Una vez cargado podremos interactuar con él para mostrar los datos en la gráfica y en la tabla de log.

\subsubsection{Información del usuario}

En esta parte de la aplicación podremos ver el nombre, foto y curso del usuario y el \emph{host} en el que se ha conectado.

\imagen{userop}{Información del usuario.}{0.9}

\subsubsection{Menú}

\imagen{menu}{Menú.}{0.9}

En este menú podremos hacer las siguientes acciones:

\begin{itemize}
	\tightlist
	\item
	\textbf{Archivo:}
	\begin{itemize}
		\tightlist
		\item
		\textbf{Cambiar asignatura:} Vuelve a la pantalla de bienvenida, que muestra la lista de cursos matriculados del usuario.
		\item
		\textbf{Guardar gráfico como...:} Abre una ventana para que el usuario guarde el gráfico en formato de imagen (.jpg o .png).
		\item
		\textbf{Cerrar sesión:} Vuelve a la pantalla de inicio de sesión.
		\item
		\textbf{Salir:} Cierra la aplicación.
	\end{itemize}
	\item
	\textbf{Editar:}
	\begin{itemize}
		\tightlist
		\item
		\textbf{Borrar selección:} Borra las selecciones de los participantes y eventos que estén seleccionados y recarga el gráfico y la tabla de log.
	\end{itemize}
	\item
	\textbf{Ayuda:}
	\begin{itemize}
		\tightlist
		\item
		\textbf{Acerca de UBULog:} Abre en el navegador un enlace al repositorio del proyecto.
	\end{itemize}
\end{itemize}













\bibliographystyle{plain}
\bibliography{bibliografiaAnexos}

\end{document}
